\section{Theoretical Foundations and Related Work}
\label{sec:background}
\subsection{Analogy involves more than mapping}
Structure-mapping theory distinguishes relational correspondence from isolated attribute overlap \citep{gentner1983structure}. A useful mapping preserves role consistency and connected relational systems. Systematicity favors interconnected, particularly causal, relations over disconnected local matches \citep{clement1991systematicity}. Computational structure mapping formalizes constraints on such alignments \citep{falkenhainer1989sme}, while alignable differences clarify which discrepancies matter for transfer \citep{gentner1997mapping}.

Access to a source and sound inference from that source are distinct. Surface resemblance can support reminding even when the retrieved source lacks the needed relational structure \citep{gentner1993roles,holyoak1987surface}. Thus, selecting or generating a plausible example does not establish that it can solve the target. Likewise, a mapping between examples may remain tied to their particular objects. Comparing solved analogues can support a schema that preserves the common solution principle and its conditions of application \citep{gick1983schema,gentner2003learning}. Adaptation then instantiates that schema under the target's own constraints; transfer may fail even when a relevant source is available \citep{novick1991math}.

These distinctions motivate our four components. They are a synthesis for evaluation, rather than a new claim that cognitive theory contains exactly four universal steps. The framework's methodological contribution is to translate them into explicit artifacts, information boundaries, and interventions. A fluent list of similarities, a general statement of a principle, and a successful target solution are assessed separately.

\subsection{Recognizing, generating, and using analogies}
Narrative benchmarks such as ARN and AnaloBench assess structural analogy identification under variations in surface similarity and context \citep{sourati2024arn,ye2024analobench}. Such tasks help establish whether a model can identify a corresponding source, but do not directly show that it can use the source solution to solve a new target. A recent synthesis connects cognitive theories of analogy to NLP research and emphasizes relational understanding beyond entity similarity \citep{petersen2026modelling}.

For problem solving, analogical prompting generates relevant exemplars before solving the target \citep{yasunaga2024analogical}. SCUA evaluates the value of teacher-generated analogies for scientific concept understanding \citep{yuan2024boosting}. These approaches establish practical benefits under their protocols, but source production, source correctness, additional explanation, and target solving can vary together. \citet{qin2025relevant} directly challenge an interpretation based on relevance alone by showing that random self-generated examples can also help.

Our work builds on these findings in two ways. We first compare analogy conditions with controls targeting specific alternative explanations, rather than treating an accuracy gain as sufficient evidence of transfer. We then independently vary structural correspondence and source-answer availability. Finally, the stage framework examines the observable intermediates and tests continuation after correcting them. It complements recognition benchmarks and prompting methods by asking where a transfer attempt becomes inadequate and what remains difficult after that failure is repaired.
